\section{PCIe: controller, DMA and Linux driver}
\label{sec:pcie}

\subsection{From the hard block to the endpoint}
The Kintex-7 has an integrated PCIe block (\code{PCIE\_2\_1}~\cite{ug477}) with eight GTX transceivers. It is
configured for Gen2 x8 with a 128-bit AXI-Stream interface at 250\,MHz (\code{userclk2}); the block's core runs at
500\,MHz (\code{userclk1}) and buffers packets in BRAMs placed next to it. LitePCIe's PHY wraps the block and moves
the streams into the sys domain through asynchronous FIFOs. On top, the LitePCIe endpoint parses and builds
transaction-layer packets (TLPs) and routes them between a memory-mapped slave (BAR0, the CSR bus) and the DMA
engines.

\paragraph{The PHY fix.} With LitePCIe 2026.08, the stock 7-series PHY took the receive byte enables of the
128-bit datapath from \code{tkeep}. On this card that added a garbage 17th beat to every 256-byte read completion,
and the DMA reader stalled after about 64\,KB. Two early host freezes happened while debugging this stall. The
project's own PHY (\file{boards/dupcb/s7pciephy\_rx\_be.py}, \code{S7PCIEPHYRXBE}) derives the byte enables from
the end-of-packet information instead. With it, LitePCIe's DMA loopback test runs at 13.43\,Gb/s
(\SI{1.68}{GB/s}, Gen2 x4) with no errors, and no host freeze has happened since, with the IOMMU on.

\subsection{DMA engines and the DDR bridge}
LitePCIe has a DMA \emph{reader} (host to card, H2C) and a \emph{writer} (card to host, C2H). Each works through a
table of descriptors (host address, length) that the driver fills; the driver runs both in \emph{loop mode} over a
ring of 256 buffers of 8\,KB in host memory, so the engines go round the ring until they are stopped. The reader
issues memory-read requests of up to 512\,B and receives completions of 256\,B; the writer sends memory-write TLPs and
raises an MSI interrupt every 32 buffers. A loop-status register tells the host which descriptor the writer is on and
how many times it has gone round.

The bridge (\file{hw/gemm/ddr\_dma.py}, \cref{fig:pcie_dma}) connects these 128-bit streams to the 512-bit
LiteDRAM native ports. One CSR selects the channel for both directions; per channel there is a
\code{LiteDRAMDMAWriter} (stream to DDR) and a \code{LiteDRAMDMAReader} (DDR to stream), each with base address,
length and enable registers. Because the reader keeps streaming the ring until it is stopped, the bridge takes exactly
\code{length} bytes and drops (and counts) everything after them or while the selected writer is not running. This
makes padding to whole DMA buffers harmless.

\begin{figure}[t]
  \centering
  \resizebox{\linewidth}{!}{% Host <-> DDR data path: loop-mode DMA rings in host memory, LitePCIe DMA engines, the DDR DMA bridge.
\begin{tikzpicture}[x=1cm,y=1cm, every node/.append style={font=\sffamily\footnotesize}]
  \tikzset{buf/.style={draw=black!60, fill=#1, minimum width=3.2mm, minimum height=5mm, inner sep=0pt},
           buf/.default=white,
           w/.style={minimum width=#1}}
  % ---- host rings ----
  \foreach \i in {0,...,5} { \node[buf=cBlue!15] (h\i) at (\i*0.36, 1.1) {}; }
  \node[font=\scriptsize] at (2.35,1.1) {\dots};
  \node[buf=cBlue!15] (h7) at (2.75,1.1) {};
  \node[note, anchor=south, text width=36mm] (nh) at (1.35,1.45) {H2C ring: 256\x8\,KB (2\,MB),\\ write-only mmap};
  \draw[->, cBlue!70!black, thick] (h7.east) to[out=0, in=0, looseness=1.6] ++(0,-0.62) -- ++(-2.85,0)
        to[out=180, in=180, looseness=1.6] (h0.west);
  \node[lbl, text=cBlue!70!black] at (1.35,0.32) {loop mode: read over and over};

  \foreach \i in {0,...,5} { \node[buf=cOrange!20] (c\i) at (\i*0.36,-2.4) {}; }
  \node[font=\scriptsize] at (2.35,-2.4) {\dots};
  \node[buf=cOrange!20] (c7) at (2.75,-2.4) {};
  \node[note, anchor=north, text width=36mm] (nc) at (1.35,-2.75) {C2H ring: 256\x8\,KB, mmap;\\ progress from the writer's\\ loop status (index, count)};
  \begin{scope}[on background layer]
    \node[region=black, fit=(h0)(h7)(c0)(c7)(nh)(nc), inner sep=5pt] (hm) {};
  \end{scope}
  \node[regionl=black, anchor=south west] at (hm.north west) {host memory (IOMMU-mapped)};

  % ---- LitePCIe DMA ----
  \node[sys, w=30mm, minimum height=13mm] (rd) at (6.3, 1.1)
       {DMA \textbf{reader} (H2C)\\[-1pt]\scriptsize descriptor table: addr, len\\[-1pt]\scriptsize MRd 512\,B \(\rightarrow\) CplD 256\,B};
  \node[sys, w=30mm, minimum height=13mm] (wr) at (6.3,-2.4)
       {DMA \textbf{writer} (C2H)\\[-1pt]\scriptsize descriptor table, loop status\\[-1pt]\scriptsize MWr TLPs, MSI per 32 buffers};
  \draw[wbus] (3.35,1.1) -- node[lbl, above] {PCIe reads} (rd.west);
  \draw[wbus] (wr.west) -- node[lbl, above] {PCIe writes} (3.35,-2.4);

  % ---- bridge ----
  \node[sys, w=17mm] (cv1) at (9.6, 1.1) {converter\\[-1pt]\scriptsize 128\(\rightarrow\)512\,b};
  \node[sys, w=13mm, trapezium, trapezium left angle=70, trapezium right angle=70, shape border rotate=270,
        minimum height=11mm, inner sep=1pt] (dmx) at (11.5, 1.1) {\scriptsize demux};
  \node[sys, w=17mm] (cv2) at (9.6,-2.4) {converter\\[-1pt]\scriptsize 512\(\rightarrow\)128\,b};
  \node[sys, w=13mm, trapezium, trapezium left angle=70, trapezium right angle=70, shape border rotate=90,
        minimum height=11mm, inner sep=1pt] (mx) at (11.5,-2.4) {\scriptsize mux};
  \node[sys, w=26mm] (w0) at (14.2, 1.75) {\code{LiteDRAMDMAWriter}\\[-1pt]\scriptsize CH0: base, length, enable};
  \node[sys, w=26mm] (w1) at (14.2, 0.45) {\code{LiteDRAMDMAWriter}\\[-1pt]\scriptsize CH1: base, length, enable};
  \node[sys, w=26mm] (r0) at (14.2,-1.75) {\code{LiteDRAMDMAReader}\\[-1pt]\scriptsize CH0: base, length, enable};
  \node[sys, w=26mm] (r1) at (14.2,-3.05) {\code{LiteDRAMDMAReader}\\[-1pt]\scriptsize CH1: base, length, enable};
  \draw[wbus] (rd) -- (cv1);
  \draw[wbus] (cv1) -- (dmx);
  \draw[wbus] (dmx.north east) ++(0,-0.15) |- (w0.west);
  \draw[wbus] (dmx.south east) ++(0,0.15) |- (w1.west);
  \draw[wbus] (r0.west) -- ++(-0.4,0) |- ([yshift=2mm]mx.east);
  \draw[wbus] (r1.west) -- ++(-0.4,0) |- ([yshift=-2mm]mx.east);
  \draw[wbus] (mx) -- (cv2);
  \draw[wbus] (cv2) -- (wr);
  \node[cdc/.style={}, draw=cRed!70!black, fill=cRed!6, font=\sffamily\scriptsize, rounded corners=1pt,
        align=left, text width=37mm, inner sep=3pt] (drop) at (11.0,-0.65)
       {\textbf{sel} CSR picks the channel for both directions. Data arriving while the selected writer is not
        running, or after \code{length} bytes, is dropped and counted (\code{h2c\_dropped}).};
  \draw[wbus] (w0.east) -- ++(0.55,0) node[right, font=\sffamily\scriptsize] {CH0 write port};
  \draw[wbus] (w1.east) -- ++(0.55,0) node[right, font=\sffamily\scriptsize] {CH1 write port};
  \draw[wbus] ([xshift=5.5mm]r0.east) node[right, font=\sffamily\scriptsize] {CH0 read port} -- (r0.east);
  \draw[wbus] ([xshift=5.5mm]r1.east) node[right, font=\sffamily\scriptsize] {CH1 read port} -- (r1.east);
  \begin{scope}[on background layer]
    \node[region=cBlue, fit=(cv1)(w0)(r1)(drop), inner sep=6pt] (brg) {};
  \end{scope}
  \node[regionl=cBlue, anchor=south west] at (brg.north west) {\code{PCIeDDRBridge} (\file{hw/gemm/ddr\_dma.py}), sys 125\,MHz};
\end{tikzpicture}
}
  \caption{The host-to-DDR data path. The driver's loop-mode rings live in host memory; the bridge converts between
    the 128-bit DMA streams and the 512-bit DDR words and selects the channel.}
  \label{fig:pcie_dma}
\end{figure}

\subsection{The Linux driver}
The kernel module is LitePCIe's driver, generated with each build (\file{<build>/driver/kernel}). It creates
\code{/dev/litepcie0} and offers:
\begin{itemize}
  \item an ioctl for 32-bit CSR reads and writes (BAR0 cannot be mapped by user space while the driver owns the
        device);
  \item ioctls that enable or disable the DMA reader and writer, which (re)program the descriptor tables from buffer~0
        in loop mode;
  \item \code{mmap} of both rings into the user process, so the host copies data straight into and out of the DMA
        buffers;
  \item MSI interrupts that count finished buffers.
\end{itemize}
The rings are allocated by the driver and mapped for the device through the IOMMU, which stays enabled. The
driver's probe resets the SoC, which is why DDR calibration has to come after the driver is loaded
(\cref{sec:flow}). A small setup script (\file{scripts/setup\_permissions.sh}) installs udev rules and a helper so
the whole bring-up runs without a password.

\subsection{Transfer protocol and throughput}
\label{sec:pcie:rates}
\code{DdrDma} (\file{sw/gemm\_host/src/ddr\_dma.cpp}) moves data in chunks of one ring (2\,MB):
\begin{description}[font=\sffamily\small, leftmargin=1.2em]
  \item[H2C] Copy the chunk into the ring, program the bridge writer (base, length, enable), enable the DMA reader,
        wait until the channel reports \code{drained} (writer done, FIFO empty), then disable the reader and the
        bridge writer. The reader cannot stay enabled between chunks, because it would keep streaming the ring.
  \item[C2H] Enable the DMA writer \emph{once} per call. For each chunk: program the bridge reader for the chunk (in
        whole 8\,KB buffers, at least two), wait until the writer's loop status shows that many more buffers done,
        copy out of the ring. Disable the writer at the end.
\end{description}

The generated driver's DMA stop waits a fixed \SI{1}{ms} (\code{udelay(1000)}), so that read completions still in
flight land before the engine is reset. With one stop per 2\,MB chunk this is about \SI{40}{\percent} of the transfer
time (\cref{fig:dma}, right). Enabling the C2H writer once per call instead of once per chunk removed it from that
direction: C2H went from 0.84 to \SI{1.49}{GB/s}, verified with random and counter patterns on both channels and
with fully verified GEMMs. H2C stays at \SI{0.82}{GB/s}: shortening the driver delay, or making the bridge
back-pressure the reader after \code{length} bytes, would remove the remaining cost, but a completion arriving after a
reset could stall the PCIe path, the same class of fault as the early host freezes, so this is left as a decision
(\cref{sec:lessons}).

\begin{figure}[t]
  \centering
  % Host <-> DDR rate per transfer size, from the GEMM runs.
\begin{tikzpicture}
\begin{axis}[rpt, height=5.8cm, xmode=log, log basis x=2, xmin=0.1, xmax=600, ymin=0, ymax=1.9,
    xtick={0.125,0.5,2,8,32,128,512}, xticklabels={128\,K,512\,K,2\,M,8\,M,32\,M,128\,M,512\,M},
    xlabel={bytes per transfer}, ylabel={GB/s}, legend style={at={(0.02,0.98)}, anchor=north west}]
  \draw[cGrey, densely dashed] (axis cs:0.1,1.679) -- (axis cs:600,1.679) node[pos=0.98, below left, font=\sffamily\tiny, text=cGrey] {raw DMA loopback 1.68};
  \addplot[l3, mark=*]        table[x=mib, y=gbps] {data/dma_c2h.csv};     \addlegendentry{C2H (C), writer enabled once per read}
  \addplot[l7, mark=o, densely dotted] table[x=mib, y=gbps] {data/dma_c2h_old.csv}; \addlegendentry{C2H before that fix}
  \addplot[l1, mark=square*]  table[x=mib, y=gbps] {data/dma_h2c.csv};     \addlegendentry{H2C (A, B)}
\end{axis}
\end{tikzpicture}
\\[2mm]
  % Time per 2 MB chunk by phase (ddr_dma_test, 64 MB, both channels).
\begin{tikzpicture}
\begin{axis}[rpt, xbar stacked, enlarge x limits=false, width=0.84\linewidth, bar width=10pt, height=4.0cm, xmin=0, xmax=2.9, xlabel={ms per 2\,MB chunk},
    ytick={0,1,2}, yticklabels={H2C, C2H before fix, C2H after fix}, y dir=reverse, enlarge y limits=0.3,
    legend style={at={(0.5,1.03)}, anchor=south, legend columns=4}, grid=none, xmajorgrids]
  \addplot[s7] table[y=x, x=copy]  {data/manual_dma_phases.csv}; \addlegendentry{ring copy}
  \addplot[s1] table[y=x, x=start] {data/manual_dma_phases.csv}; \addlegendentry{start}
  \addplot[s3] table[y=x, x=wait]  {data/manual_dma_phases.csv}; \addlegendentry{data moving}
  \addplot[s4] table[y=x, x=stop]  {data/manual_dma_phases.csv}; \addlegendentry{stop (driver: 1\,ms delay)}
\end{axis}
\end{tikzpicture}

  \caption{Host \(\leftrightarrow\) DDR transfers. Top: rate per transfer size, from every GEMM run (A and B
    uploads, C downloads). Small transfers pay the per-chunk start and stop; above \(\sim\)8\,MB the rate is flat.
    Bottom: time per 2\,MB chunk by phase, from \code{ddr\_dma\_test}. The red part is the driver's fixed 1\,ms stop,
    gone from C2H after the fix.}
  \label{fig:dma}
\end{figure}
